\section{USAMO 1985}
        \begin{enumerate}
        	\item (\textit{USAMO 1985}) Determine whether or not there are any positive integral solutions of the simultaneous equations 

 \[x_1^2+x_2^2+\cdots+x_{1985}^2=y^3, \hspace{20pt} x_1^3+x_2^3+\cdots+x_{1985}^3=z^2\]
with distinct integers $x_1,x_2,\cdots,x_{1985}$.


 



\textbf{Solution}

Lemma: For a positive integer $n$, $1^3+2^3+\cdots +n^3 = (1+2+\cdots +n)^2$ (Also known as Nicomachus's theorem)

 Proof by induction: The identity holds for $1$. Suppose the identity holds for a number $n$. It is well known that the sum of first $n$ positive integers is $\frac{n(n+1)}{2} = \frac{n^2+n}{2}$. Thus its square is $\frac{n^4+2n^3+n^2}{4}$. Adding $(n+1)^3=n^3+3n^2+3n+1$ to this we get $\frac{n^4+6n^3+13n^2+12n+4}{4}$, which can be rewritten as $\frac{(n^4+4n^3+6n^2+4n+1)+2(n^3+3n^2+3n+1)+(n^2+2n+1)}{4}$ This simplifies to $\frac{(n+1)^4+2(n+1)^3+(n+1)^2}{4} = ({\frac{(n+1)^2+(n+1)}{2}})^2 = (1+2+\cdots +n+(n+1))^2$. The induction is complete.

 
Let $j$ be the sum $1+2+\cdots 1985$, and let $k$ be the sum $1^2 + 2^2 + \cdots + 1985^2$. Then assign $x_i$ the value $ik^4$ for each $i = 1, 2,\cdots 1985$. Then:
 \begin{align*} x_1^2 +x_2^2 +\cdots +x_{1985}^2 & = 1^2k^8 +2^2k^8+\cdots +1985^2k^8 = k^8(1^2+2^2+\cdots +1985^2) = k^9 = {(k^3)}^3\\ x_1^3 +x_2^3 +\cdots +x_{1985}^3 & = 1^3k^{12}+2^3k^{12}+\cdots 1985^3k^{12}=k^{12}(1^3+2^3+\cdots 1985^3) = k^{12}j^2 = ({k^6j})^2 \end{align*}


 Thus, a positive integral solution exists.

 -Circling

 


	\item (\textit{USAMO 1985}) Determine each real root of


 $x^4-(2\cdot10^{10}+1)x^2-x+10^{20}+10^{10}-1=0$


 correct to four decimal places.


 



\textbf{Solution}

The equation can be re-written as 
 \begin{align}\label{eqn1} (x+10^5)^2(x-10^5)^2 -(x+10^5)(x-10^5) -x-1=0. \end{align}

 We first prove that the equation has no negative roots.Let $x\le 0.$ The equation above can be further re-arranged as
 \begin{align*}[(x+10^5)(x-10^5)+1][(x+10^5)(x-10^5)-2]=x-1.\end{align*}
The right hand side of the equation is negative. Therefore 
 \[[(x+10^5)(x-10^5)+1][(x+10^5)(x-10^5)-2)]<0,\] and we have$-1<(x+10^5)(x-10^5) <2.$ Then the left hand side of the equation is bounded by
 \[|[(x+10^5)(x-10^5)+1][(x+10^5)(x-10^5)-2]|\le 3\times 3.\]However, since $|(x+10^5)(x-10^5)|\le 2$ and $x<0,$ it follows that $|x+10^5| <\frac{2}{|x-10^5|}<2\times 10^{-5}$ for negative $x.$ Then $x<2\times 10^{-5}-10^5.$ The right hand side of the equation is then a large negative number. It cannot be equal to the left hand side which is bounded by 9.

 Now let $x>0.$ When $x=10^5,$ the left hand side of equation (1) is negative. Therefore the equation has real roots on both side of $10^5$, as its leading coefficient is positive. We will prove that $x=10^5$ is a good approximation of the roots (within $10^{-2}$). In fact, we can solve the "quadratic" equation (1) for $(x+10^5)(x-10^5)$:
 \[(x+10^5)(x-10^5)=\frac{1\pm\sqrt{1+4(x+1)}}{2}.\]Then 
 \[x-10^5=\frac{1\pm\sqrt{1+4(x+1)}}{2(x+10^5)}.\]Easy to see that$|x-10^5| <1$ for positve $x.$ Therefore, $10^5-1<x<10^5+1.$ Then
 \begin{align*} |x-10^5|&=\left|\frac{1\pm\sqrt{1+4(x+1)}}{2(x+10^5)}\right |\\ &\le \left |\frac{1}{2(x+10^5)}\right |+\left |\frac{\sqrt{1+4(x+1)}}{2(x+10^5)}\right |\\ &\le \frac{1}{2(10^5-1+10^5)} +\frac{\sqrt{1+4(10^5+1+1)}}{2(10^5-1+10^5)} \\ &<10^{-2}. \end{align*}


 Let $x_1$ be a root of the equation with $x_1<10^5.$ Then $0<10^5-x_1<10^{-2}$ and 
 \[x_1-10^5=\frac{1-\sqrt{1+4(x_1+1)}}{2(x_1+10^5)}.\]An aproximation of $x_1$ is defined as follows:
 \[\tilde{x}_1=10^5+\frac{1-\sqrt{1+4(10^5+1)}}{2(10^5+10^5)}.\]We check the error of the estimate:
 \begin{align*} |\tilde{x}_1-x_1|&=\left | \frac{1-\sqrt{1+4(10^5+1)}}{2(10^5+10^5)}- \frac{1-\sqrt{1+4(x_1+1)}}{2(x_1+10^5)} \right | \\ &\le \left |\frac{1}{2(10^5+10^5)}- \frac{1}{2(x_1+10^5)}\right |+\left |\frac{\sqrt{1+4(10^5+1)}}{2(10^5+10^5)}- \frac{\sqrt{1+4(x_1+1)}}{2(x_1+10^5)}\right |. \end{align*}


 The first absolute value 
 \[\left |\frac{1}{2(10^5+10^5)}- \frac{1}{2(x_1+10^5)}\right | =\frac{|x_1- 10^5|}{2(10^5+10^5)(x_1+10^5)}<10^{-12}.\]

 The second absolute value 
 \begin{align*} &\left |\frac{\sqrt{1+4(10^5+1)}}{2(10^5+10^5)} - \frac{\sqrt{1+4(x_1+1)}}{2(x_1+10^5)} \right |\\ &\le \left |\frac{\sqrt{1+4(10^5+1)}}{2(10^5+10^5)}- \frac{\sqrt{1+4(x_1+1)}}{2(10^5+10^5)}\right |+\left |\frac{\sqrt{1+4(x_1+1)}}{2(10^5+10^5)}- \frac{\sqrt{1+4(x_1+1)}}{2(x_1+10^5)}\right |\\ &\le 10^{-7}+10^{-9}, \end{align*}
through a rationalized numerator.Therefore $|\tilde{x}_1-x_1|\le 10^{-6}.$

 For a real root $x_2$ with $x_2>10^5,$ we choose 
 \[\tilde{x}_2=10^5+\frac{1+\sqrt{1+4(10^5+1)}}{2(10^5+10^5)}.\]We can similarly prove it has the desired approximation.

 \textbf{Notes} Another round of iteration can increase the accuracy to more than 10 decimal places:
 \begin{align*} \tilde{x}_1^\prime=10^5+\frac{1-\sqrt{1+4(\tilde{x}_1+1)}}{2(\tilde{x}_1+10^5)},\\ \tilde{x}_2^\prime=10^5+\frac{1+\sqrt{1+4(\tilde{x}_2+1)}}{2(\tilde{x}_2+10^5)}. \end{align*}


 J.Z.

 


	\item (\textit{USAMO 1985}) Let $A,B,C,D$ denote four points in space such that at most one of the distances $AB,AC,AD,BC,BD,CD$ is greater than $1$. Determine the maximum value of the sum of the six distances.


 



\textbf{Solution}

Suppose that $AB$ is the length that is more than $1$. Let spheres with radius $1$ around $A$ and $B$ be $S_A$ and $S_B$. $C$ and $D$ must be in the intersection of these spheres, and they must be on the circle created by the intersection to maximize the distance. We have $AC + BC + AD + BD = 4$.

 In fact, $CD$ must be a diameter of the circle. This maximizes the five lengths $AC$, $BC$, $AD$, $BD$, and $CD$. Thus, quadrilateral $ACBD$ is a rhombus. 

 Suppose that $\angle CAD = 2\theta$. Then, $AB + CD = 2\sin{\theta} + 2\cos{\theta}$. To maximize this, we must maximize $\sin{\theta} + \cos{\theta}$ on the range $0^{\circ}$ to $90^{\circ}$. However, note that we really only have to solve this problem on the range $0^{\circ}$ to $45^{\circ}$, since $\theta > 45$ is just a symmetrical function.

 For $\theta < 45$, $\sin{\theta} \leq \cos{\theta}$. We know that the derivative of $\sin{\theta}$ is $\cos{\theta}$, and the derivative of $\cos{\theta}$ is $-\sin{\theta}$. Thus, the derivative of $\sin{\theta} + \cos{\theta}$ is $\cos{\theta} - \sin{\theta}$, which is nonnegative between $0^{\circ}$ and $45^{\circ}$. Thus, we can conclude that this is an increasing function on this range. 

 It must be true that $2\sin{\theta} \leq 1$, so $\theta \leq 30^{\circ}$. But, because $\sin{\theta} + \cos{\theta}$ is increasing, it is maximized at $\theta = 30^{\circ}$. Thus, $AB = \sqrt{3}$, $CD = 1$, and our sum is $5 + \sqrt{3}$.

 


	\item (\textit{USAMO 1985}) There are $n$ people at a party. Prove that there are two people such that, of the remaining $n-2$ people, there are at least $\lfloor n/2\rfloor -1$ of them, each of whom knows both or else  knows neither of the two. Assume that "know" is a symmetrical relation; $\lfloor x\rfloor$ denotes the greatest integer less than or equal to $x$.


 



\textbf{Solution}

Consider the number of pairs (X, {Y, Z}), where X, Y, Z are distinct points such that X is joined to just one of Y, Z. If X is joined to just k points, then there are just k(n - 1 - k) $\leq$ (n - 1)2/4 such pairs (X, {Y, Z}). Hence in total there are at most $\frac{n(n - 1)^2}{4}$ such pairs. But there are $\frac{n(n - 1)}{2}$ possible pairs {Y, Z}. So we must be able to find one of them {A, B} which belongs to at most $\lfloor \frac{n-1}{2} \rfloor$ such pairs. 

 Hence there are at least $n - 2 - \lfloor \frac{n - 1}{2} \rfloor = \lfloor \frac{n}{2} \rfloor - 1$ points X which are joined to both of A and B or to neither of A and B. (If confused by the floor, consider n = 2m and n = 2m+1 separately!)

 


	\item (\textit{USAMO 1985}) Let $a_1,a_2,a_3,\cdots$ be a non-decreasing sequence of positive integers. For $m\ge1$, define $b_m=\min\{n: a_n \ge m\}$, that is, $b_m$ is the minimum value of $n$ such that $a_n\ge m$. If $a_{19}=85$, determine the maximum value of
$a_1+a_2+\cdots+a_{19}+b_1+b_2+\cdots+b_{85}$.


 



\textbf{Solution}

We create an array of dots like so: the array shall go out infinitely to the right and downwards, and at the top of the $i$th column we fill the first $a_i$ cells with one dot each. Then the $19$th row shall have 85 dots. Now consider the first 19 columns of this array, and consider the first 85 rows. In row $j$, we see that the number of blank cells is equal to $b_j-1$. Therefore the number of filled cells in the first 19 columns of row $j$ is equal to $20-b_j$.

 We now count the number of cells in the first 19 columns of our array, but we do it in two different ways. First, we can sum the number of dots in each column: this is simply $a_1+\cdots+a_{19}$. Alternatively, we can sum the number of dots in each row: this is $(20-b_1)+\cdots +(20-b_{85})$. Since we have counted the same number in two different ways, these two sums must be equal. Therefore
 \[a_1+\cdots +a_{19}+b_1+\cdots +b_{85}=20\cdot 85=\boxed{1700}.\]Note that this shows that the value of the desired sum is constant.

 


\end{enumerate}